\chapter{Toolbar System and Navigation}
\label{ch:toolbar}

\emph{Pull requests: \#111, \#123, \#126, \#132 (toolbar part), \#262.}

\section{Overview}
With node and grillage overlays in place (Chapter~\ref{ch:cad}), the next need was a clean way
to \emph{drive} the 3D view: fit, zoom, pan, rotate, zoom-to-window, and toggle the various
overlays. Rather than scatter this logic across the CAD widget, I introduced a dedicated
\textbf{toolbar controller} that owns every toolbar button and wires it to the correct
operation on whichever view is active---the 3D CAD view or the 2D plots view. This chapter
covers the controller, the label-visibility toggles built on top of it, and the later
state-management fixes.

\section{The Toolbar Controller}
The core of this work is a new module,
\texttt{desktop/ui/utils/toolbar\_controller.py}, containing the
\texttt{ToolBarController} class (PR~\#111). It separates the toolbar's behaviour from the CAD
widget: the controller finds the toolbar buttons, connects them to camera operations, and can
\texttt{reset()} itself to tear down all connections when the active view changes.

\begin{lstlisting}[caption={Camera operations wired by the controller (\texttt{toolbar\_controller.py}).},label={lst:tbctrl}]
class ToolBarController:
    def _cad_zoom_fit(self):
        """Fit all visible geometry in the OCC view."""
        try:
            if cad_widget.display is not None and not cad_widget._cad_init_pending:
                cad_widget.display.FitAll()
        except Exception:
            pass

    def _cad_zoom_in(self):   """Zoom in by a fixed 10% step."""
    def _cad_zoom_out(self):  """Zoom out by a fixed 10% step."""
    def _cad_toggle_pan(self):          ...   # delegate to viewer pan mode
    def _cad_toggle_zoom_window(self):  ...   # delegate to viewer zoom-window
    def _cad_toggle_rotate(self):       ...   # delegate to viewer rotate mode
\end{lstlisting}

Seven camera operations are exposed in total: \emph{zoom fit}, \emph{zoom in}, \emph{zoom
out}, \emph{pan}, \emph{zoom window}, \emph{rotate}, and \emph{grillage toggle}. The pan,
zoom-window, and rotate toggles delegate to the navigation modes I had added to the OCC viewer
(\texttt{set\_navigation\_mode} in \texttt{custom\_3dviewer.py}, see
Chapter~\ref{ch:cad})---for example the zoom-window mode uses the OCC \texttt{ZoomArea()} call
with a \texttt{WindowFit()} fallback.

The controller binds to whichever view is active through two public methods, and cleans up
after itself:
\begin{lstlisting}[caption={Binding the controller to a view, and resetting it.},label={lst:tbbind}]
def bind_to_cad_3d(self, cad_widget):
    """Wire the toolbar buttons to the 3D CAD view."""

def bind_to_plots(self, plots_widget):
    """Wire the toolbar buttons to the 2D plots view."""

def reset(self):
    """Remove all connections and restore the plain button appearance."""
\end{lstlisting}

\section{Bug Fixes and Cleanup}
PR~\#123 stabilised the first iteration of the toolbar. It corrected the button
\emph{highlight} states (buttons were showing as active when they were not), removed the
temporary \emph{Zoom Fit / Zoom In / Zoom Out} buttons that had been used during development,
made the pan operation work on the plot view, and fixed bugs in the navigation cube
(\emph{navcube}) behaviour. This was the point at which the toolbar moved from a working
prototype to a clean, production-ready control surface.

\section{Label-Visibility Toggles}
PR~\#126 added a group of toggle buttons for the on-model labels: \emph{node numbers},
\emph{element numbers}, and \emph{girder labels}, available in both the 3D CAD view and the 2D
plots. The toolbar buttons themselves are created as checkable buttons
(\texttt{custom\_widgets.py}), the controller locates them by tooltip and wires them, and the
plot widget applies the visibility to the underlying Matplotlib artists.

\begin{lstlisting}[caption={Creating the checkable label-toggle buttons (\texttt{custom\_widgets.py}).},label={lst:tbtoggles}]
self.btn_show_node_number = create_button(
    ":/vectors/tool_bar/Node_Number.svg", "Node Number")
self.btn_show_node_number.setCheckable(True)

self.btn_show_element_number = create_button(
    ":/vectors/tool_bar/Element_number.svg", "Element Number")
self.btn_show_element_number.setCheckable(True)

self.btn_show_girder_label = create_button(
    ":/vectors/tool_bar/Girder_label.svg", "Girder Labels")
self.btn_show_girder_label.setCheckable(True)
\end{lstlisting}

In the plots view, toggling a button simply flips a flag and re-applies visibility to the
tagged text artists, then requests a redraw---a data-driven approach that avoids rebuilding the
plot:

\begin{lstlisting}[caption={Applying element-number visibility to Matplotlib artists (\texttt{mpl\_plot\_widget.py}).},label={lst:mplvis}]
def _on_element_numbers_toggled(self, checked: bool):
    self._show_element_numbers = checked
    self._apply_element_number_visibility()
    self._canvas.draw_idle()

def _apply_element_number_visibility(self):
    for ax in self._fig.axes:
        for text in ax.texts:
            if text.get_gid() == "element_number":
                text.set_visible(self._show_element_numbers)
\end{lstlisting}

\section{State Management and Final Polish}
The last round of toolbar work (PR~\#262) addressed three usability issues that surfaced once
the toolbar was in daily use.

\subsection{Resetting toggles on unlock}
When the model is cleared or the inputs are unlocked, every CAD toggle should return to its
off state so the UI does not lie about what is currently shown. I generalised the reset so it
iterates over \emph{all} managed buttons---Grillage, Node, Node Number, Element Number, Rotate,
Pan, Zoom Window, Axis, and Legend---rather than a hard-coded subset:

\begin{lstlisting}[caption={Reset every managed toggle when the CAD overlays clear (\texttt{toolbar\_controller.py}).},label={lst:tbreset}]
def sync_cad_overlays_off(self):
    for btn in self._managed_buttons:
        self._sync_btn_to(btn, False)
\end{lstlisting}

\subsection{Disabling the design button after a design}
After running a design, the inputs freeze (lock). Previously the \emph{Design} button remained
enabled, so a user could re-trigger a design with the exact same, now-frozen inputs. The fix
ties the button's enabled state directly to the lock state inside
\texttt{input\_dock.py}:

\begin{lstlisting}[caption={Disable the design button whenever the inputs are locked (\texttt{input\_dock.py}).},label={lst:designbtn}]
def toggle_lock(self):
    ...
    self.design_btn.setDisabled(self.is_locked)
\end{lstlisting}

A matching disabled style was added to the custom button stylesheet so the state is visually
obvious:

\begin{lstlisting}[language=,caption={Disabled-button styling (\texttt{custom\_buttons.py}).},label={lst:disabledqss}]
QPushButton:disabled {
    background-color: #D6E3A3;
}
\end{lstlisting}

\subsection{Hand cursor for pan}
Finally, to give the pan tool the feedback users expect, the viewer now shows an open-hand
cursor while pan mode is active and the normal arrow otherwise, driven from
\texttt{set\_navigation\_mode} in \texttt{custom\_3dviewer.py}.

\section{Result}
OsdagBridge gained a single, coherent toolbar that drives both the 3D CAD and 2D plot views:
fit, zoom, pan, rotate, and zoom-window navigation; independent toggles for nodes, node
numbers, element numbers, girder labels, the grillage, the axis, and the legend; consistent
button highlight states; and correct enable/disable behaviour tied to the design lifecycle.

\osdagfig{figure4.1.png}{The OsdagBridge toolbar, driving the 3D CAD and plot views.}{fig:toolbar}

\osdagfig{figure4.2.png}{Node-number, element-number, and girder-label toggles
applied to the model.}{fig:toolbarlabels}

\section{Summary}
The toolbar controller centralised all view-navigation logic into one reusable class, the
label toggles gave the user fine-grained control over on-model annotations, and the final
state-management fixes made the toolbar behave correctly across the full design lifecycle. The
\emph{Axis} toggle referenced above is the subject of Chapter~\ref{ch:axis}.
